\section{总结与延伸}

\subsection{一条清晰的技术主线}

回顾陈天奇近 20 年的职业生涯，一条主线贯穿始终：\textbf{让机器学习系统更高效、更易用、更可部署}。每一个项目都是对前一个项目瓶颈的回应：

\begin{itemize}
    \item 做 KDD Cup 发现需要高效的树模型 $\rightarrow$ \textbf{XGBoost}
    \item 做 XGBoost 时重新燃起对深度学习的热情 $\rightarrow$ \textbf{CXXNet/MXNet}
    \item 做 MXNet 发现工程人力是瓶颈 $\rightarrow$ \textbf{TVM}（用编译技术替代手写 kernel）
    \item 做 TVM 需要在垂直领域验证 $\rightarrow$ \textbf{MLC LLM}（大模型端侧部署）
    \item TVM 的商业化需求 $\rightarrow$ \textbf{OctoML}（开源 $\rightarrow$ 创业 $\rightarrow$ 被 NVIDIA 收购）
\end{itemize}

\begin{table}[H]
\centering
\begin{tabular}{p{2.2cm}p{1.8cm}p{9cm}}
\toprule
\textbf{项目} & \textbf{时期} & \textbf{解决的问题} \\
\midrule
XGBoost & 2014 & 让树模型在大规模数据上跑得更快、更易用（自动处理 missing value） \\
MXNet & 2015 & 让深度学习框架支持 Python first、自动调度、多卡并行、自动求导 \\
TVM & 2017 & 让 ML 模型自动适配各种硬件（GPU/TPU/NPU），减少手写 kernel \\
MLC LLM & 2023 & 让大语言模型跑在手机、浏览器、MacBook 等端侧设备上 \\
\bottomrule
\end{tabular}
\caption{陈天奇的技术贡献时间线}
\end{table}

这条主线的底层逻辑是：\textbf{ML 系统不应该那么复杂，应该让更少的人通过正确的技术就能快速搭建}。这个信念从 XGBoost 时代（一个人做到极致）到 TVM 时代（编译器自动生成 kernel）到 MLC LLM 时代（端侧部署跨平台），一以贯之。

\subsection{对年轻研究者的启示}

综合访谈内容，陈天奇给年轻人的核心建议可以总结为以下几点，每一点都有他亲身经历的支撑：

\begin{enumerate}
    \item \textbf{挑正确的问题}：``做平庸的问题和重要的问题花的时间几乎一样。''这是交大敖平老师教给他的，也是他 20 年来反复验证的。ImageNet 是正确的问题，RBM 是错误的方法------问题的选择比方法的选择更重要。

    \item \textbf{不要同时锁死问题和方法}：这是他在交大做 ImageNet 两年无果后学到的``最宝贵的教训''。挑了问题后，方法要灵活调整------后来在 KDD Cup 中果断从 RBM 转向矩阵分解，就是这个教训的应用。

    \item \textbf{做到极致}：Carlos 的要求是``write the best paper in the world''。不是追求论文数量，而是每一个项目都要做到 best paper level。XGBoost 的成功就是``一个算法做到极致''的最好例证。

    \item \textbf{接受失败}：``做很难的问题必然会意味着不确定性和失败。''早期失败反而是好事------让你更敢放手做事，避免路径依赖。

    \item \textbf{追随内心}：选择时问自己``如果失败了会不会后悔''。这不是理性的投入产出分析，而是确保走的是理想中的路。

    \item \textbf{从 first principle 出发}：不怕进入陌生领域。从写 transpiler 到做编译器到设计 NPU，陈天奇每次都是``出生牛犊不怕虎''地直接上手。系统设计能力可以在实践中磨练出来。

    \item \textbf{开源是高校学生的独特机遇}：在公司里你从大项目的小部件做起；在开源项目中你有机会从头架构一个系统。这种训练在其他地方很难获得。
\end{enumerate}

\begin{knowledgebox}{给 PhD 学生的特别建议}
当被问到``如果你今天是三年级 PhD 学生，只能选择修复或改变一件事情，会做什么''时，陈天奇的回答是：``我应该会选择做我现在在做的事情。''也就是继续做机器学习编译------让更少的人用更少的 engineering resource 解决 ML 系统在最新硬件上的 bring up 问题。他不会因为时代不同就改变方向，因为这个问题``远远没有被解决''。
\end{knowledgebox}

\subsection{对 AI 未来的展望}

陈天奇对 AGI 的定义不感兴趣（``我对 AGI 本身怎么定义并不是特别感兴趣''），更关心``pragmatically 有用''：能不能让一个人的团队撬动 100 人团队的产出？现在的技术``的确能够让我们更快地向那个目标前进''------比如他自己在用 Cursor 做 TVM 重构，如果没有 AI 工具``做重构会困难很多''。

对于 AI 世界的未来格局，他希望看到百花齐放------开源和个体都能做创新，而不是几家巨头主导一切。当被问到``那样的世界出现的概率是在增加还是减少''时，他的回答很谨慎：``我只能说它需要每一个人的努力。肯定会有不同人会有不同的信念。''

他坚持的技术路线始终是：\textbf{用更少的资源、更自动地做更多的事情}。``或许它在几年之内会被解决，或许可能不一定会被解决，但至少它是一个很有趣的问题。''

这句话也许是对陈天奇整个职业生涯最好的概括：不追求确定性的成功，而是追求有趣的问题。只要问题足够重要、足够有趣，就值得投入------即使结果是失败，也是``再出发的过程''。

\subsection{结语}

访谈的最后，主持人留了一段开放时间让陈天奇自由发挥。他说：

\textbf{``在整个过程中我非常幸运能够获得众多人的帮助。最让我感觉收获深刻的是：我们还是可以去接受失败。在接受失败的过程中，继续鼓励自己做有趣的事。每个过程其实都是一个再出发的过程。We need to reinvent ourselves, then keep doing the next thing。''}

如果让他给过去和未来的自己各说一句话，他的选择不是从现在回头教过去的自己什么，而是相反------他需要过去那个``出生牛犊不怕虎''的自己来提醒现在的自己：\textbf{记住对自己的承诺，坚持理想，往下走下去。}

主持人总结这段访谈说：理想的 10 年后 20 年后的陈天奇，应该是一个``有更多次失败的陈天奇------并且能够继续保持那一个最初的心态去挑战我们想要挑战的东西''。陈天奇欣然同意：``会有更多的失败，并且可以继续保持那个最初的心态去挑战。这样子最好了。''

\begin{importantbox}{访谈核心金句}
以下是本次访谈中最具启发性的几段原话：
\begin{itemize}
    \item ``你去做一件比较平庸的问题和一个重要的问题，花的时间几乎是一样的。''
    \item ``时间是有限的，所以一定要做想要做的事情。''
    \item ``很多时候你的勇气可能是你过去的自己给你自己的约定。''
    \item ``如果选了这条路并且失败了，你还会不会后悔？''
    \item ``失败也没那么可怕。然后要接受可以失败这个事实。''
    \item ``我觉得现在的我最难的事情，还是要保持初心。''
    \item ``做优秀的技术始终是没有错的。''
    \item ``We need to reinvent ourselves, then keep doing the next thing.''
\end{itemize}
\end{importantbox}

\subsection{拓展阅读}

\begin{itemize}
    \item 陈天奇，《机器学习科研的十年》，知乎，2019。这篇文章回顾了他从 2009 年到 2019 年的科研历程，其中很多故事在本次访谈中有更详细的展开。文章中``随机过程总会收敛到和付出相称的稳态''这句话被广泛传播。
    \item XGBoost 论文：Chen \& Guestrin, ``XGBoost: A Scalable Tree Boosting System'', KDD 2016。这篇论文至今引用近 7 万次，是 KDD 历史上引用最高的论文之一。
    \item MXNet 论文：Chen et al., ``MXNet: A Flexible and Efficient Machine Learning Library for Heterogeneous Distributed Systems'', NeurIPS Workshop 2015。作者按字母序排列，体现了 PhD 学生之间``纯粹的合作''精神。
    \item TVM 论文：Chen et al., ``TVM: An Automated End-to-End Optimizing Compiler for Deep Learning'', OSDI 2018。开创了机器学习编译这个研究方向。
    \item LIME 论文：Ribeiro, Singh, Guestrin, ``Why Should I Trust You? Explaining the Predictions of Any Classifier'', KDD 2016。Carlos 实验室同年另一篇开创性工作，开启了可解释机器学习方向。
    \item MLC LLM 项目：\url{https://github.com/mlc-ai/mlc-llm}。将大语言模型编译部署到各种端侧设备。
    \item Apache TVM 项目：\url{https://tvm.apache.org/}。机器学习编译的开源基础设施。
    \item MLC 课程（陈天奇在 CMU 讲授的机器学习编译课程）：\url{https://mlc.ai/}。系统讲解机器学习编译的原理和实践。
    \item XGrammar：\url{https://github.com/mlc-ai/xgrammar}。结构化数据生成的开源库，已成为 LLM Agent 结构化输出的事实标准之一。
    \item WhynotTV Podcast 第二期（胡渊鸣专访）：同样涉及开源项目商业化、学术创业等话题，与本期形成有趣的对照。
\end{itemize}

\end{document}
